\section{Related Work}
\label{sec:related}

\paragraph{Imitation learning and robotic demonstrations.}
Behavior cloning directly regresses expert actions and has a long history in
autonomous control \citep{pomerleau1989alvinn}. Surveys of learning from
demonstration emphasize the importance of representation, demonstration
quality, and correspondence between teacher and learner
\citep{argall2009survey,osa2018survey}. Offline BC avoids environment
interaction during training but suffers covariate shift. DAgger addresses this
by collecting expert labels on learner-induced states
\citep{ross2011dagger}; our study deliberately retains a fixed offline dataset
without interactive data collection. Offline manipulation studies show
sensitivity to demonstration quality, algorithm design, and stopping criteria
\citep{mandlekar2022what}; data-quality work formalizes curation through
distribution shift \citep{belkhale2023data}. We instead vary staged sampling
over one fixed scripted-expert dataset.

\paragraph{Curriculum learning for sequential decision-making.}
Curriculum learning orders training examples from easier to harder and can
improve optimization or generalization in some settings
\citep{bengio2009curriculum}. Automated curriculum research broadens this
view to task selection and teacher--student formulations
\citep{portelas2020survey}. Reverse curriculum generation has also shown that
starting near successful states and expanding outward can aid sparse-reward
control \citep{florensa2017reverse}. Our Reverse Learner condition is more
limited: it reverses a fixed demonstration ranking and does not generate new
tasks or states. Cross-Episodic Curriculum (CEC) instead structures sequences
of online learning trials and mixed-quality demonstrations within a
Transformer's context to improve agent learning across episodes
\citep{shi2023cross}. Recent preprint work, CurriBuffer, applies
difficulty-aware demonstration selection, curriculum sampling, and retrieval
to Vision-Language-Action policy specialization
\citep{hosseini2026curribuffer}. These approaches establish the
relevance of organizing experience for sequential and robotic policies. Our
setting is deliberately narrower: it uses fixed offline BC data and compares
difficulty definitions without cross-episodic context or retrieval.
Classification studies further show that scoring and pacing are distinct
design choices \citep{hacohen2019power}, and that curriculum,
anti-curriculum, and random schedules depend on training budget and data noise
\citep{wu2021when}.

\paragraph{Difficulty estimation and curriculum design.}
Example difficulty may be based on task metadata, teacher effort, uncertainty,
loss, or forgetting. These signals describe different phenomena, and
loss-based signals risk in-sample bias. Rampp et al. systematically study the
similarity and robustness of difficulty scoring functions and their effects on
curriculum learning in image and acoustic-scene classification
\citep{rampp2024difficulty}. We extend this question to trajectory-level BC and
closed-loop robotic manipulation. Specifically, we compare Task difficulty
(initial geometry), Expert difficulty (demonstration effort), and Learner
difficulty (five-fold out-of-fold prediction error). These are descriptive
signals rather than ground-truth measures of complexity.

Taken together, CEC and CurriBuffer motivate curriculum organization for
sequential and robotic learning, while Rampp et al. motivate comparing how
difficulty is defined. We hold the demonstrations, policy family, and training
budget fixed while comparing three trajectory rankings and five staged
sampling schedules under physical closed-loop evaluation.
